\documentclass[12pt]{article}
\usepackage[amsmath]{e-jc}
\usepackage{mathtools}
\usepackage[T1]{fontenc}
\usepackage{lmodern}
\usepackage{microtype}
\hypersetup{
  pdftitle={Path Ramsey Thresholds in Coprime Graphs},
  pdfauthor={Alexander Richard Towell},
  pdfsubject={AI-assisted editorial draft; human mathematical review pending; not submitted},
  pdfkeywords={coprime graph, path Ramsey number, extremal graph theory}
}

\newcommand{\Pcop}{P_{\mathrm{cop}}}
\newcommand{\Fzero}{\mathcal F_0}
\newcommand{\Ic}{\overline K}
\newcommand{\repo}{https://github.com/queelius/computational-explorations}
\newcommand{\revision}{8bda81e8d3f867bb1450b0e492592164416ba98f}

% The journal style is used for preparation, not to imply acceptance.
\makeatletter
\def\the@dateline{Editorial revision, October 20, 2026\\[0.5ex]
Human mathematical review pending; not submitted}
\def\ps@plain{%
  \def\@oddhead{}\def\@evenhead{}%
  \def\@oddfoot{\normalfont\hfil Working draft --- \thepage\hfil}%
  \let\@evenfoot\@oddfoot}
\makeatother
\pagestyle{plain}
\MSC{05C38, 05C55, 05D10, 11A05}
\Copyright{The author.}
\title{Path Ramsey Thresholds in Coprime Graphs}
\author{Alexander Richard Towell\authornote{1}}
\authortext{1}{Department of Computer Science, Southern Illinois University
Edwardsville, Edwardsville, Illinois, USA
(\email{lex@metafunctor.com}). ORCID: 0000-0001-6443-9897.}

\begin{document}
\maketitle

\begin{abstract}
Let $G_n$ be the graph on $\{1,\ldots,n\}$ in which two distinct
vertices are adjacent when they are coprime. Write $\Pcop(k)$ for
the least $n$ such that every two-coloring of $E(G_n)$ contains a
monochromatic path with $k$ edges. We prove that $\Pcop(9)=17$.
An elementary coloring gives the lower bound. For the upper bound,
explicit path certificates force any coloring avoiding a nine-edge
path to have a connected, 48-edge color graph on 17 vertices in
which every four-vertex deletion leaves at least five edges.
We classify the possible graphs and show that each admits a
four-vertex deletion leaving at most one edge. We also prove
\[
 2\leq\liminf_{k\to\infty}\frac{\Pcop(k)}k
 \leq\limsup_{k\to\infty}\frac{\Pcop(k)}k
 \leq\frac{197\pi^2}{595},
\]
and give independently checkable certificates for the previously
reported values $\Pcop(7)=\Pcop(8)=13$.
\end{abstract}

\smallskip\noindent\textbf{Keywords:} coprime graph, path Ramsey number,
extremal graph theory, computational certificate.

\section{Introduction}

How large must the coprime graph be before every red--blue edge
coloring contains a monochromatic path of a prescribed length?
Let $G_n$ have vertex set $[n]=\{1,\ldots,n\}$ and edge set
\[
 E(G_n)=\bigl\{\{a,b\}:1\leq a<b\leq n,\ \gcd(a,b)=1\bigr\}.
\]
We write $\Pcop(k)$ for the least $n$ that forces a monochromatic
path with $k$ edges. Throughout the paper, $P_j$ denotes a path
on $j$ \emph{vertices}; thus the target defining $\Pcop(k)$ is
$P_{k+1}$. All paths are simple. Each color graph has the full
vertex set of the host, including any isolated vertices.

Parity gives a useful starting point. The even vertices of $G_n$
are independent, so the odd vertices form a vertex cover.
Partitioning that cover into two small sets gives an edge coloring
in which both colors have small vertex covers and therefore no
long paths. On $G_{16}$, this construction avoids a monochromatic
nine-edge path. Our main result shows that it cannot be extended
to $G_{17}$.

\begin{theorem}\label{thm:main}
Every two-coloring of $E(G_{17})$ contains a monochromatic $P_{10}$,
whereas some two-coloring of $E(G_{16})$ contains none. Consequently,
$\Pcop(9)=17$.
\end{theorem}

The upper bound combines three finite properties of $G_{17}$ with
an extremal graph argument. The finite properties are certified
by explicit paths. They reduce a hypothetical avoiding coloring
to two connected color graphs with 47 and 48 edges, each retaining
at least five edges after any four vertices are deleted.
The following theorem gives the opposing conclusion for the
48-edge graph.

We use $X\vee Y$ for the join of disjoint graphs $X$ and $Y$, and
$\Ic_t$ for an independent set of order~$t$.

\begin{theorem}\label{thm:classification}
Let $J$ be a connected $P_{10}$-free graph on 17 vertices with
48 edges. Either $J$ has a vertex cover of size at most four, or
\[
 J\cong K_3\vee(K_3\cup\Ic_{11}).
\]
In either case, some four-vertex deletion leaves at most one edge.
\end{theorem}

The proof of Theorem~\ref{thm:classification} uses the stability
theorems of F\"uredi, Kostochka, and Verstra\"ete~\cite{fkv} and
F\"uredi, Kostochka, Luo, and Verstra\"ete~\cite{fkvl}, together
with the longest-path classification of Ali and Staton~\cite{ali}
and Hamiltonian closure. The graph lies exactly at the edge-count
threshold of the relevant stability result, so a separate
analysis of equality is needed.

For general path lengths, the same vertex-cover construction
gives a lower bound. An application of the Erd\H{o}s--Gallai
theorem to a suitable induced subgraph gives an upper bound.

\begin{theorem}\label{thm:linear}
For $k\geq3$,
\[
 \Pcop(k)\geq4\left\lfloor\frac{k-1}{2}\right\rfloor+1.
\]
Moreover,
\[
 2\leq\liminf_{k\to\infty}\frac{\Pcop(k)}k
 \leq\limsup_{k\to\infty}\frac{\Pcop(k)}k
 \leq\frac{197\pi^2}{595}=3.267751373\ldots.
\]
In particular, $\Pcop(k)=\Theta(k)$.
\end{theorem}

The path problem differs from the corresponding problem for
cliques. Du, Xi, Deng, and Ma~\cite{du} express the edge-coloring
thresholds for complete targets in terms of classical Ramsey
numbers and prime labels. For paths, pulling a coloring back
through such labels can identify distinct vertices, so the
same argument does not preserve the required path.
Their Table~12 records earlier path computations from the
author's repository, and their Open Problem~9.2 asks about
non-complete edge targets.

Section~\ref{sec:finite} gives the Ramsey argument for
Theorem~\ref{thm:main}, using Theorem~\ref{thm:classification}.
We prove the classification in Section~\ref{sec:classification}
and the linear bounds in Section~\ref{sec:linear}.
Section~\ref{sec:earlier} records certificates for the earlier
values $\Pcop(7)=\Pcop(8)=13$. The final section discusses
the remaining questions, and Appendix~\ref{sec:evidence}
gives the computational evidence and checking instructions.

\section{Proof of the exact threshold}\label{sec:finite}

\subsection{A coloring on sixteen vertices}

If a graph has a vertex cover of size $m$, every path has at most
$2m$ edges: it uses at most $m$ vertices of the cover and at most
$m+1$ vertices outside it. We will use this observation both here
and in the general lower bound.

In $G_{16}$, partition the odd vertices into
\[
 A=\{1,3,5,7\},\qquad B=\{9,11,13,15\}.
\]
Color an edge red if it meets $A$, and blue otherwise.
Then $A$ covers the red graph. Since there are no even--even
edges, every blue edge meets $B$, so $B$ covers the blue graph.
Both colors therefore have paths of at most eight edges.
This proves $\Pcop(9)>16$.

\subsection{Three properties of the host}

For disjoint sets $A,B\subseteq[17]$, let $G_{17}[A,B]$
denote the graph consisting of the host edges between $A$ and $B$.
We need the following three properties. They will force large
color components, prevent sparse remainders after vertex
deletions, and rule out the final obstruction to connectedness.

\begin{lemma}\label{lem:finite}
The graph $G_{17}$ has the following properties.
\begin{enumerate}
\item\label{item:cuts} If $A\cup B=[17]$ is a partition with
$|A|,|B|\geq5$, then $G_{17}[A,B]$ contains a $P_{10}$.
\item\label{item:faults} For every four-element set $A\subset[17]$
and every set $F$ of at most four edges of $G_{17}-A$, the graph
$G_{17}-A-F$ contains a $P_{10}$.
\item\label{item:extensions} If $|A|=4$, $e(G_{17}[A])\geq5$,
and $B=[17]\setminus A$, then, for every edge
$e\in E(G_{17}[B])$, the graph with edge set
$E(G_{17}[A,B])\cup\{e\}$ contains a $P_{10}$.
\end{enumerate}
\end{lemma}

\begin{proof}[Proof by finite certificates]
For parts~\ref{item:cuts} and~\ref{item:extensions}, the archive
supplies a path for every required cut or pair $(A,e)$.
For part~\ref{item:faults}, it supplies, for each four-vertex
deletion up to checked host symmetries, a family of paths that
cannot all be met by four edges. Hence at least one path
survives any further deletion of at most four edges.

An independent checker verifies every path, the coverage of all
required cases, and the symmetries used to reduce the vertex
deletions. The path-family condition is checked by exhaustive
branching over possible sets of four edges.
Appendix~\ref{sec:evidence} describes this verification and
identifies the certificates and checker.
\end{proof}

\subsection{From the host properties to connected color graphs}

We first recall the exceptional graph in the connected-path
bounds. The graph $H(n,k,a)$ has three parts $A,B,C$ of sizes
$a,n-k+a,k-2a$, respectively. It has all edges within $A\cup C$
and all edges between $A$ and $B$, and no others. Thus
\begin{equation}\label{eq:h}
 h(n,k,a):=e(H(n,k,a))
 =\binom{k-a}{2}+a(n-k+a).
\end{equation}
In particular, deleting the four-vertex part $A$ from
$H(s,10,4)$ leaves one edge and isolated vertices.

For a connected 13-vertex $P_{10}$-free graph,
Balister, Gy\H{o}ri, Lehel, and Schelp~\cite[Theorem~1.3]{bgls}
give
\begin{equation}\label{eq:13bound}
 e\leq\max\left\{\binom82+5,\binom52+4\cdot8\right\}=42.
\end{equation}
For a connected $P_{10}$-free graph of order $s\geq14$,
the path stability theorem~\cite[Theorem~1.6]{fkv} gives
$e\leq3s-3$, unless the graph is a subgraph of $H(s,10,4)$.

\begin{proposition}\label{prop:reduction}
If a two-coloring of $G_{17}$ has no monochromatic $P_{10}$,
both color graphs are connected on all 17 vertices, their edge
counts are 47 and 48, and each retains at least five edges after
every four-vertex deletion.
\end{proposition}

\begin{proof}
We first establish the deletion condition. Suppose that deleting
a four-element set $A$ leaves at most four red edges. Delete
those edges as well. Every remaining edge is blue, so
Lemma~\ref{lem:finite}(\ref{item:faults}) gives a blue $P_{10}$,
a contradiction. The same argument applies with the colors
interchanged.

Next, each color must have a component of order at least 13.
Otherwise, a union of its components would have order between
5 and 12. Indeed, if no single component has order in that range,
all components have order at most four; accumulate them until
the union first has at least five vertices, when it has at most
eight. All host edges across the resulting cut have the other
color, contrary to Lemma~\ref{lem:finite}(\ref{item:cuts}).

Let $s\geq13$ be the order of the large component in either
color. If $s\geq14$, this component cannot be a subgraph of
$H(s,10,4)$: deleting its distinguished four vertices would
leave at most
\[
 1+\binom{17-s}{2}\leq4
\]
edges of that color in the whole graph, contradicting the
deletion condition. The connected-path bounds therefore give
the following upper bounds for the total number of edges in
that color:
\[
\begin{array}{c|rrrrr}
 \text{large component order }s&13&14&15&16&17\\ \hline
 \text{total edges, at most}&48&42&43&45&48.
\end{array}
\]
The first entry is $42+\binom42$; the others are
$3s-3+\binom{17-s}{2}$. Since $e(G_{17})=95$, neither
color can have fewer than 47 edges. Thus the two edge counts
are 47 and 48, and each large component has order 13 or 17.

It remains to exclude order 13. Suppose a component of one
color has vertex set $B$, with $|B|=13$, and put
$A=[17]\setminus B$. By~\eqref{eq:13bound}, at least
$47-42=5$ edges of that color lie in $A$. In particular,
$e(G_{17}[A])\geq5$. Every crossing edge has the other color.
If any host edge inside $B$ also had that other color,
Lemma~\ref{lem:finite}(\ref{item:extensions}) would give a
monochromatic $P_{10}$. Consequently every host edge inside
$B$ has the component's color. But
Lemma~\ref{lem:finite}(\ref{item:faults}), with no edge
deletion, gives a $P_{10}$ in $G_{17}[B]$.
This contradiction proves that both color graphs are connected
on all 17 vertices.
\end{proof}

\begin{proof}[Proof of Theorem~\ref{thm:main}]
The coloring above proves the lower bound. For the upper bound,
suppose that a coloring of $G_{17}$ avoids a monochromatic
$P_{10}$. Proposition~\ref{prop:reduction} supplies a connected,
48-edge color graph in which every four-vertex deletion leaves
at least five edges. Theorem~\ref{thm:classification} supplies
a four-vertex deletion leaving at most one edge, a contradiction.
\end{proof}

\section{The extremal equality case}\label{sec:classification}

We now prove Theorem~\ref{thm:classification}. The argument has
three stages. First, a dense bipartite configuration in the
universal augmentation of $J$ forces a four-vertex cover.
If this configuration is absent, the stability theorems let us
delete vertices of degree three until only ten vertices remain.
We then classify the possible ten-vertex cores and determine how
they can be reconstructed without creating a $P_{10}$.

\subsection{Stability inputs and a forbidden configuration}

Adding a universal vertex $u$ to a connected $P_{10}$-free graph
$L$ of order at least two gives a 2-connected graph $L+u$
with no cycle of length at
least 11. Indeed, a cycle of that length through $u$ would
leave a $P_{10}$ after $u$ is removed, while a cycle avoiding
$u$ would itself contain a $P_{10}$. We call $L+u$ the
\emph{universal augmentation} of $L$.

Let $\Fzero$ be the family obtained from $K_{5,6}$ by deleting
at most two edges. For a member of this family, write $A$ and
$B$ for its parts of sizes five and six.
The following specializations of the published stability
results will be used. In each statement, $Q$ is 2-connected,
has order $n$, and has circumference $c(Q)<11$.
\begin{enumerate}
\item If $n\geq15$ and $e(Q)>h(n,11,4)$, then
$Q\subseteq H(n,11,5)$ \cite[Theorem~1.4]{fkv}.
\item If $n\geq11$ and
$e(Q)>\max\{h(n,11,4),h(n,11,3)\}$, then
$Q\subseteq H(n,11,5)$ or $Q\subseteq H(n,11,2)$
\cite[Theorem~2.3]{fkvl}.
\item If $n\geq11$, a subgraph of $H(n,11,5)$ with more
than $h(n,11,4)$ edges contains a member of $\Fzero$
\cite[Lemma~4.6]{fkvl}.
\item If an edge contraction $Q/xy$ is 2-connected and
contains a member of $\Fzero$, then $Q$ contains a member
of $\Fzero$ \cite[Lemma~4.9]{fkv}.
\end{enumerate}
Contraction suppresses loops and parallel edges. The relevant
edge counts from~\eqref{eq:h} are
\begin{equation}\label{eq:thresholds}
 h(n,11,4)=4n-7,\qquad h(n,11,3)=3n+4,\qquad
 h(n,11,2)=2n+18.
\end{equation}

We also use the following path property of $\Fzero$
\cite[Lemma~4.8]{fkv}. Every pair of distinct vertices in a
member of $\Fzero$ is joined by a path of at least eight edges.
If the pair is not entirely in $A$, there is such a path of at
least nine edges. A path between two vertices of $B$ has even
length, so in that case there is a path of ten edges.

\begin{lemma}\label{lem:F0}
Suppose that a 2-connected graph $Q$ has circumference less than
11 and contains a member of $\Fzero$ with parts $A,B$.
Then $Q-A$ is independent. If $Q=J+u$, where $u$ is universal
and $|J|\geq10$, then $J$ has a vertex cover of size at most four.
\end{lemma}

\begin{proof}
Let $S$ be a component of $Q-(A\cup B)$. It has at least two
attachments in $A\cup B$. If $|S|\geq2$, the attachment edges
also meet at least two vertices of $S$; otherwise their unique
endpoint in $S$ would be a cut vertex. Thus we can choose
attachment edges $px,qy$ with distinct $p,q\in A\cup B$
and distinct $x,y\in S$. A path in $S$ from $x$ to $y$
gives a $p$--$q$ path of at least three edges, internally
outside $A\cup B$. Combining it with a path of at least eight edges
in the member of $\Fzero$ gives a cycle of length at least
11, a contradiction. Hence every such component $S$ is
a singleton.

If a singleton has an attachment in $B$, join that attachment
to a second attachment by a path of at least nine edges in
$\Fzero$. Together with the two attachment edges, this again
gives a cycle of length at least 11. Therefore all vertices
outside $A\cup B$ attach only to $A$.
An edge within $B$ would close a ten-edge path in $\Fzero$
and is likewise impossible. This proves that $Q-A$ is
independent.

In a universal augmentation, $u$ has degree at least ten.
Every vertex outside $A$ has degree at most five, so $u\in A$.
The set $A\setminus\{u\}$ is consequently a four-vertex cover
of $J$.
\end{proof}

For the rest of the proof, assume that $J$ has no vertex cover
of size at most four. By Lemma~\ref{lem:F0}, its universal
augmentation contains no member of $\Fzero$. The same is
true of the augmentation of every induced subgraph of $J$.

\subsection{Deleting vertices down to a ten-vertex core}

We first show that every edge of $J$ lies in at least two
triangles. The augmentation $J+u$ has 18 vertices and 65 edges.
If an edge $xy$ lies in at most one triangle of $J$, it lies
in at most two triangles of $J+u$. Contracting it therefore
loses at most three edges, leaving at least 62 edges on
17 vertices. The contracted graph is still a universal
augmentation of a connected graph, hence is 2-connected.
Contraction cannot increase circumference.

Since $62>h(17,11,4)=61$, the first and third stability
inputs give a member of $\Fzero$ after contraction.
The fourth input lifts it back to $J+u$, a contradiction.
In particular, every vertex of $J$ has degree at least three.

We will preserve the following conditions for a current induced
subgraph $L$ of order $n$:
\begin{equation}\label{eq:invariant}
 \begin{gathered}
 L\text{ is connected and }P_{10}\text{-free},\qquad e(L)=3n-3,\\
 \delta(L)\geq2,\qquad L+u\text{ contains no member of }\Fzero.
 \end{gathered}
\end{equation}
They hold initially for $L=J$. Suppose that $n\geq11$.

As before, an edge lying in at most one triangle of $L$
can be contracted in $L+u$ at a loss of at most three edges.
The augmentation has $n+1$ vertices and $4n-3$ edges,
so the contraction has $n$ vertices and at least $4n-6$
edges. This exceeds both $4n-7$ and $3n+4$.
The second stability input therefore places the contraction
in $H(n,11,5)$ or $H(n,11,2)$. The first alternative is
impossible by the dense-subgraph and lifting inputs.
For $n\geq13$, the second is also impossible because
$4n-6>2n+18=h(n,11,2)$.

For $n=11,12$, we exclude the second alternative by counting
leaves. Put $M=L/xy$. An old vertex can become a leaf only
if it is a common neighbor of $x$ and $y$, and there is at
most one such vertex. The merged vertex can become a leaf
only if $x$ and $y$ both had degree two and shared their other
neighbor. If this common neighbor also became a leaf, the
three vertices would form an isolated triangle in $L$,
contrary to connectedness and $n>3$. Thus $M$ has at most
one leaf.

On the other hand, if $M+u\subseteq H(n,11,2)$, the universal
vertex $u$ must lie in the two-vertex joining part of the
container. Removing $u$ would make $M$ a connected spanning
subgraph of
\[
 K_1\vee(K_7\cup\Ic_{n-9}).
\]
Each of its $n-9\geq2$ independent vertices must then be a
leaf, a contradiction. We have proved that every edge of
$L$ lies in at least two triangles. Hence $\delta(L)\geq3$.

We next exclude $\delta(L)\geq4$. If a longest path in a
connected graph is not spanning, it has at least $2\delta+1$
vertices, where $\delta$ is the minimum degree. Here is the
short argument. Suppose the path is $v_1,\ldots,v_r$ with
$r\leq2\delta$. All neighbors of its endpoints lie on the
path. Among the $r-1$ indices $i$, at least $\delta$ satisfy
$v_1v_{i+1}\in E(L)$ and at least $\delta$ satisfy
$v_iv_r\in E(L)$. The two index sets intersect, closing the
path into a cycle through all its vertices. Connectivity and
a vertex outside the cycle then give a longer path, a
contradiction.

In our graph $L$, a longest path has order $r\leq9<n$.
Thus $\delta(L)\geq4$ would force $\delta(L)=4$ and $r=9$.
Ali and Staton~\cite[Theorem~3]{ali} classify connected graphs
of minimum degree $\delta\geq2$, longest path order
$2\delta+1$, and order at least $2\delta+2$.
They are subgraphs of $K_\delta\vee\Ic_{n-\delta}$ or
windmills $K_1\vee(tK_\delta)$.
For $\delta=4$, a windmill has $5(n-1)/2<3n-3$ edges.
In the other case, minimum degree four forces every edge
between the parts of $K_4\vee\Ic_{n-4}$ to be present.
Its augmentation then contains $K_{5,6}$, also impossible.
Therefore $\delta(L)=3$.

Let $v$ be a degree-three vertex. Each edge incident with $v$
lies in two triangles, so the three neighbors of $v$ form
a triangle. Deleting $v$ preserves connectedness, removes
exactly three edges, and leaves minimum degree at least two.
The other conditions in~\eqref{eq:invariant} are inherited.
We may therefore repeat the deletion.

Seven deletions leave a connected $P_{10}$-free graph $L_{10}$
on ten vertices with 27 edges and minimum degree at least two.
Its predecessor $L_{11}$ had minimum degree at least three.
Reversing the process adds one vertex at a time, adjacent to
exactly three mutually adjacent vertices already present.

\subsection{The possible cores}

We show that $L_{10}$ has one of the following three forms:
\begin{equation}\label{eq:cores}
 \begin{aligned}
 &K_2\vee(K_5\cup\Ic_3),\\
 &K_3\vee(K_3\cup\Ic_4),\\
 &\text{a subgraph of }K_4\vee\Ic_6.
 \end{aligned}
\end{equation}
Let $Q=L_{10}+u$. It has 11 vertices, 37 edges, and minimum
degree at least three, and it has no Hamiltonian cycle.
Take its 11-closure: repeatedly add an edge between nonadjacent
vertices whose degree sum is at least 11.
The Bondy--Chv\'atal theorem~\cite{closure} preserves
non-Hamiltonicity.

If any edge is added, the closure has at least 38 edges.
The two thresholds in the second stability input both equal
37 at $n=11$, so the closure is a subgraph of $H(11,11,5)$
or $H(11,11,2)$. The latter has two vertices of degree two
and is excluded by the minimum degree.
In $H(11,11,5)=K_5\vee\Ic_6$, the universal vertex $u$
must lie in the five-vertex part. Removing it gives
$L_{10}\subseteq K_4\vee\Ic_6$.

It remains to consider the case where $Q$ is already closed.
Write its degrees in nondecreasing order as
$d_1,\ldots,d_{11}$. Since $Q$ is not Hamiltonian,
Chv\'atal's degree-sequence criterion, in the form stated
in~\cite[Theorem~4.2]{fkvl}, gives an index $i<11/2$ with
\[
 d_i\leq i,\qquad d_{11-i}\leq10-i.
\]
Minimum degree three restricts $i$ to $3,4,5$.

For $i=3$, deleting $u$ gives the following termwise upper
bounds for the degrees of $L_{10}$:
\[
 (2,2,2,6,6,6,6,6,9,9).
\]
Their sum is $54=2e(L_{10})$, so equality holds in every
position. The two degree-nine vertices are universal.
Removing them leaves five vertices of degree four and three
isolated vertices, which forces
$L_{10}=K_2\vee(K_5\cup\Ic_3)$.

For $i=4$, the corresponding upper sequence is
\[
 (3,3,3,3,5,5,5,9,9,9),
\]
again with sum 54. The three degree-nine vertices are
universal; removing them leaves a triangle and four isolated
vertices. Hence $L_{10}=K_3\vee(K_3\cup\Ic_4)$.

For $i=5$, choose a set $B$ of six vertices of degree at most
five in $Q$, and let $A=V(Q)\setminus B$. We will show that
an edge inside $B$ would force $Q$ to be Hamiltonian.
First, $A$ is a clique. The degree sum of $Q$ is 74, only
six below the upper bound
\[
 6\cdot5+5\cdot10=80.
\]
If two vertices of $A$ were nonadjacent, closedness would
give them degree sum at most ten. Those two vertices alone
would contribute a deficit of at least ten, a contradiction.

Let $b=e(Q[B])$, and let $m$ be the number of missing
$A$--$B$ edges. Since $A$ is a clique, the edge count gives
$37=10+30-m+b$, or $m=b+3$.
For $v\in B$, let $s_v$ count its missing cross edges and
$t_v$ its neighbors within $B$. Its degree is
$5-s_v+t_v\leq5$. Thus, on putting $q_v=s_v-t_v$, we obtain
\[
 q_v\geq0,\qquad
 \sum_{v\in B}q_v=m-2b=3-b,\qquad 0\leq b\leq3.
\]
For $a\in A$, let $r_a$ count its missing cross edges.
Every absent pair $av$ satisfies, by closedness,
\begin{equation}\label{eq:closed-pair}
 (10-r_a)+(5-q_v)\leq10,
 \qquad\text{so}\qquad r_a+q_v\geq5.
\end{equation}

Suppose $b>0$. Since $q_v\leq3-b$, every vertex of $A$
incident with a missing cross edge must have at least
$b+2$ such missing edges. There are only $m=b+3$ missing
cross edges in total. Consequently they are all incident
with a single vertex $a_0\in A$.
It follows that $s_v\leq1$ for every $v\in B$.
As $t_v\leq s_v$, the edges in $B$ form a matching.

If $b=1$, the endpoints of its edge have $q_v=0$.
Equation~\eqref{eq:closed-pair} would require $r_{a_0}\geq5$,
whereas $m=4$. If $b=2$, the missing cross edges join $a_0$
to the four matching endpoints and one isolated vertex.
If $b=3$, they join $a_0$ to all six vertices of $B$.
In both remaining cases,
\[
 Q\cong K_4\vee(3K_2\cup K_1).
\]
Indeed, the four vertices of $A\setminus\{a_0\}$ are universal.
Outside them there are three edges and one isolated vertex;
when $b=2$, the third edge joins $a_0$ to its sole neighbor
in $B$.

This graph is Hamiltonian: place its three two-vertex blocks
and its singleton between successive vertices of the $K_4$
in a cycle. Thus $b>0$ is impossible. We have $b=0$,
and deleting the universal vertex $u\in A$ again gives
$L_{10}\subseteq K_4\vee\Ic_6$.
This completes the list~\eqref{eq:cores}.

\subsection{Reconstructing the graph}

We now examine which cores can survive the seven reverse
additions. Each new vertex has a triangle as its neighborhood
among the vertices already present.

The first core, $K_2\vee(K_5\cup\Ic_3)$, is impossible.
Its three independent vertices have degree two, so all three
would have to be adjacent to the first new vertex in order
to obtain $\delta(L_{11})\geq3$. But three independent
vertices cannot form its required neighborhood triangle.

For the second core, $K_3\vee(K_3\cup\Ic_4)$, let $A$
be the joining triangle, $C$ the other triangle, and $B$
the independent set. Every vertex outside $A$ is an endpoint
of a nine-vertex path, as the following templates show:
\[
 \begin{split}
 &c_1,c_2,c_3,a_1,b_1,a_2,b_2,a_3,b_3,\\
 &b_1,a_1,c_1,c_2,c_3,a_2,b_2,a_3,b_3.
 \end{split}
\]
A new vertex cannot be adjacent to such an endpoint, since
prepending the new vertex would give a $P_{10}$.
Its earlier neighborhood must therefore be exactly $A$.
The same templates apply to this new vertex as a member of
the enlarged independent set $B$. Induction forces all
seven additions to have earlier neighborhood $A$.
No edges join the added vertices, and the final graph is
$K_3\vee(K_3\cup\Ic_{11})$.

Finally, suppose $L_{10}\subseteq K_4\vee\Ic_6$, with parts
$A$ and $B$. The container has 30 edges and the core has 27,
so at most three cross edges are missing.
We claim that every vertex of $B$, and every previously added
vertex whose earlier neighborhood lies in $A$, is an endpoint
of a nine-vertex path.

To prove the claim, let $z$ be such a vertex. Choose a neighbor
$a_1\in A$ of $z$ and order the other three vertices as
$a_2,a_3,a_4$. Select five original vertices of $B$,
excluding $z$ if $z\in B$, and call this set $B'$.
We seek distinct representatives in $B'$ for the four slots
\[
 \{a_1,a_2\},\qquad \{a_2,a_3\},\qquad
 \{a_3,a_4\},\qquad \{a_4\}.
\]
A vertex is eligible for a slot if it is adjacent to every
vertex in that slot.

Each slot has at least two eligible vertices, since there
are five candidates and at most three missing cross edges.
This proves Hall's condition for one or two slots.
For a collection of three or four slots, a vertex outside
the union of their eligible sets must account for at least
two missing cross edges: one missing edge affects at most
two slots. With only three missing cross edges in total,
at most one vertex of $B'$ can be outside such a union.
The union therefore has at least four vertices.
Hall's condition holds for every collection of slots.

The resulting distinct representatives give the path
\[
 z,a_1,b_1,a_2,b_2,a_3,b_3,a_4,b_4,
\]
proving the claim. Now consider the first added vertex whose
earlier neighborhood is not contained in $A$. Its earlier
neighbor outside $A$ starts a path of this form, so
prepending the new vertex creates a $P_{10}$.
No such addition is possible. Thus $A$ covers the final
graph, contrary to our standing assumption.

Only the triangle-join graph remains when there is no
four-vertex cover. In that graph, deleting the joining
triangle and one vertex of the other triangle leaves one
edge. In the cover case, deleting four vertices containing
the cover leaves no edges. This proves
Theorem~\ref{thm:classification}.

\section{Linear bounds}\label{sec:linear}

\begin{proof}[Proof of Theorem~\ref{thm:linear}]
\emph{The lower bound.}
Put $m=\lfloor(k-1)/2\rfloor$. The $2m$ odd vertices cover
$G_{4m}$. Partition them into two sets of size $m$, color
every edge meeting the first set red, and color the remaining
edges blue. Each color has a cover of size $m$, so its paths
have at most $2m<k$ edges. This gives
\[
 \Pcop(k)\geq4m+1
 =4\left\lfloor\frac{k-1}{2}\right\rfloor+1
\]
and hence the asymptotic lower constant two.

\medskip
\noindent\emph{A density upper bound.}
By the Erd\H{o}s--Gallai path theorem~\cite{eg}, a graph
with no $k$-edge path has at most $(k-1)v/2$ edges.
If a two-coloring of a subgraph $H\subseteq G_n$ avoids
such a path, applying the theorem to both color graphs gives
$e(H)\leq(k-1)v(H)$. Therefore
\begin{equation}\label{eq:density}
 e(H)>(k-1)v(H)\quad\Longrightarrow\quad\Pcop(k)\leq n.
\end{equation}
The strict inequality in this criterion is essential.

Write $E_n=e(G_n)$. Counting ordered coprime pairs by
M\"obius inversion gives
\[
 2E_n+1=\sum_{r=1}^n\mu(r)\lfloor n/r\rfloor^2
       =\frac6{\pi^2}n^2+O(n\log n).
\]
The extra one counts the pair $(1,1)$.
Replacing the floors by their arguments contributes
$O(n\log n)$, and the tail of
$\sum_{r\geq1}\mu(r)/r^2=1/\zeta(2)$ contributes $O(n)$.
Taking $H=G_n$ in~\eqref{eq:density} already yields
\[
 \limsup_{k\to\infty}\frac{\Pcop(k)}k\leq\frac{\pi^2}{3}.
\]

\medskip
\noindent\emph{Improving the edge-to-vertex ratio.}
To improve the constant, we remove a set of vertices that
contributes relatively few edges. We use the multiples of
$30$, $42$, and $70$. Their small prime factors restrict
their coprime neighbors, and all these vertices are even,
so they form an independent set. Put
\[
 D_n=\{a\leq n:30\mid a\text{ or }42\mid a\text{ or }70\mid a\},
 \qquad H_n=G_n-D_n.
\]
Every pair of the three divisors has least common multiple
210. Inclusion-exclusion gives
\[
 \begin{split}
 |D_n|
 &=\lfloor n/30\rfloor+\lfloor n/42\rfloor
   +\lfloor n/70\rfloor-2\lfloor n/210\rfloor\\
 &=\frac{13}{210}n+O(1).
 \end{split}
\]

We also need to count the edges lost. For a fixed squarefree
integer $d\geq2$, let $C_d(n)$ count ordered coprime pairs
$(a,b)\in[n]^2$ with $d\mid a$. M\"obius inversion gives
\begin{align*}
 C_d(n)
 &=\sum_{r=1}^n\mu(r)
   \left\lfloor\frac n{\operatorname{lcm}(d,r)}\right\rfloor
   \left\lfloor\frac nr\right\rfloor\\
 &=\frac{6n^2}{\pi^2\prod_{p\mid d}(p+1)}
   +O_d(n\log n).
\end{align*}
The floor and tail estimates are as above. The leading
coefficient follows from the Euler product
\[
 \begin{split}
 \sum_{r\geq1}\frac{\mu(r)}{r\operatorname{lcm}(d,r)}
 &=\frac1d\prod_{p\mid d}(1-1/p)
             \prod_{p\nmid d}(1-1/p^2)\\
 &=\frac6{\pi^2\prod_{p\mid d}(p+1)}.
 \end{split}
\]
Because $D_n$ is independent, each lost edge is counted
exactly once by an ordered pair whose first endpoint is
deleted. The number of lost edges is therefore
\begin{align*}
 &C_{30}(n)+C_{42}(n)+C_{70}(n)-2C_{210}(n)\\
 &\quad=\frac6{\pi^2}
  \left(\frac1{72}+\frac1{96}+\frac1{144}-\frac2{576}\right)n^2
  +O(n\log n)\\
 &\quad=\frac{n^2}{6\pi^2}+O(n\log n).
\end{align*}
It follows that
\[
 v(H_n)=\frac{197}{210}n+O(1),\qquad
 e(H_n)=\frac{17}{6\pi^2}n^2+O(n\log n),
\]
and hence
\[
 \frac{e(H_n)}{v(H_n)}
 =\frac{595}{197\pi^2}n+O(\log n).
\]
For any $\varepsilon>0$, choose
\[
 n=\left\lceil\left(\frac{197\pi^2}{595}+\varepsilon\right)k
       \right\rceil.
\]
The positive linear margin from $\varepsilon k$ dominates
the logarithmic error, so~\eqref{eq:density} applies for
all sufficiently large $k$. Letting $\varepsilon$ tend to
zero proves the asserted upper bound.
\end{proof}

\section{Certificates for the earlier thresholds}\label{sec:earlier}

The values
\[
 \Pcop(7)=\Pcop(8)=13
\]
were reported in the author's earlier computations and appear
in Table~12 of Du et al.~\cite{du}. We record their certificates
here for completeness. The upper-bound certificate uses a
different method from the explicit host paths at $n=17$.

For the lower bound, partition the odd vertices of $G_{12}$
into $\{1,3,11\}$ and $\{5,7,9\}$. The vertex-cover coloring
used in Section~\ref{sec:finite} gives each color a cover of
size three, so neither color has a path with more than six
edges. Thus $\Pcop(7),\Pcop(8)>12$.

For the upper bound at $n=13$, assign one Boolean variable
to each of the 57 host edges. An eight-edge path gives two
clauses, excluding its two monochromatic assignments.
The archived formula contains 71\,832 valid path-exclusion
clauses and one harmless tautology ensuring that the solver
registers all 57 variables. It imposes no symmetry restrictions.

An independent semantic check verifies that every path clause
excludes a monochromatic assignment of a valid eight-edge path.
A separate DRAT certificate establishes that the formula is
unsatisfiable. This subset of the path constraints is enough:
every coloring with no monochromatic eight-edge path would
satisfy all of them. Hence every coloring of $G_{13}$ contains
an eight-edge path, and therefore also a seven-edge path.
Together with the lower coloring, this proves the two
previously reported values.

Appendix~\ref{sec:evidence} identifies the formula, DRAT proof,
verification records, and required checkers.

\section{Related bounds and further questions}\label{sec:context}

The equality case in Theorem~\ref{thm:classification} lies
just outside the strict stability bounds used in its proof.
The universal augmentation has 18 vertices and 65 edges,
exactly $h(18,11,4)$.
The later refinements of Ma and Ning~\cite{ma-ning} also
retain a strict inequality at this point.
In their notation, $W_{n,a,c}=H(n,c+1,a)$, and at
$(n,c)=(18,10)$ the term $f(n,4,c)$ equals 65.
Their Theorems~1.7 and~1.9 require more edges than a maximum
containing this term, so neither directly gives the boundary
classification needed here. This comparison concerns those
statements; a comprehensive determination of priority for
the classification is outside the scope of the paper.

The lower construction suggests asking whether
\[
 \Pcop(2m+1)=4m+1\qquad\text{for every }m\geq1.
\]
Theorem~\ref{thm:main} proves the case $m=4$, but its
host-specific certificates do not establish the general
upper bound. It also remains to close the gap between the
asymptotic constants $2$ and $197\pi^2/595$. The deletion
set used in Section~\ref{sec:linear} gives an explicit
improvement over the full-host density bound; we do not
claim that this choice is optimal or that the refinement
has priority over all earlier work.

A shorter structural proof of
$\Pcop(7)=\Pcop(8)=13$ would be useful as well.
It would give an account of these earlier values without
the separate SAT certificate.

\appendix
\section{Certificates and reproducibility}\label{sec:evidence}

The computational evidence and proof notes are available in
the repository at the immutable
\href{\repo/tree/\revision}{revision \texttt{8bda81e}}
\cite{repository}. This revision predates the present
editorial revision. All paths below are relative to the
repository root.

\subsection{What each archive contains}

\begin{description}
\item[The three host properties.]
The
\href{\repo/tree/\revision/docs/audits/coprime-path-nine-2026-10-18}
{structural archive}
contains the crossing paths, deletion families, and cut
extensions certifying the three parts of Lemma~\ref{lem:finite}.
\item[The nine-edge threshold.]
The
\href{\repo/tree/\revision/docs/audits/coprime-path-nine-resolution-2026-10-18}
{resolution archive}
contains \texttt{lower-n16-k9.json}, which records all 79
colored edges of the lower witness, together with source
and verification records.
\item[The earlier thresholds.]
The
\href{\repo/tree/\revision/docs/audits/coprime-paths-2026-10-18}
{earlier path archive}
contains the 45-edge lower witness at $n=12$, the compressed
formula and DRAT proof at $n=13$, and their verification records.
\end{description}

The first two archives use checkers written in standard-library
Python; verifying their certificates requires no SAT solver.

\subsection{Verifying the three host properties}

For Lemma~\ref{lem:finite}(\ref{item:cuts}), a certificate
gives a crossing path for each smaller side of a cut. There
are
\[
 \sum_{s=5}^8\binom{17}{s}=62\,322
\]
such sides; taking complements covers the other orientations.
For part~\ref{item:extensions}, there is one path for every
qualifying pair $(A,e)$.

For part~\ref{item:faults}, there are
$\binom{17}{4}=2380$ vertex deletions.
Checked host automorphisms reduce them to 262 equivalence
classes. For each class, the certificate supplies a family
of surviving paths with no edge transversal of size at most
four. Here an \emph{edge transversal} is a set of edges
meeting every path in the family. If no such transversal
exists, deleting four edges cannot destroy all the paths.

The checker establishes nonexistence of a transversal by
exhaustive branching. It selects a path not yet met by the
chosen edges and branches on each of that path's edges as
the next choice. Any transversal extending the current
choice must contain at least one of these edges, so the
branches cover every possibility. Branches requiring more
than four edges are rejected.

For every supplied path, the checker verifies ten distinct
vertices, all nine gcd adjacencies, and the relevant cut
or deletion restrictions. It also verifies every transposition
within each declared vertex class against host adjacency.
These transpositions generate the within-class permutations,
so deletions with the same class counts are equivalent.
The checker derives the required cases from the task
parameters and requires every deletion signature exactly
once. It does not rely on saved completion flags.

The expected counts are:
\begin{center}
\begin{tabular}{ll}
\hline
Certificate & Verified cases or witnesses\\
\hline
Cuts & $62\,322$ cuts\\
Vertex and edge deletions & $2380$ deletions in $262$ classes\\
                         & $30\,671$ paths in the families\\
Extensions & $63\,859$ pairs $(A,e)$ on $1253$ sets $A$\\
\hline
\end{tabular}
\end{center}

To obtain the cited revision and run the checks with Python
3.10 or newer:

\noindent\begin{minipage}{\linewidth}
\small
\begin{verbatim}
git clone https://github.com/queelius/computational-explorations.git
cd computational-explorations
git checkout --detach 8bda81e8d3f867bb1450b0e492592164416ba98f

cert_dir=docs/audits/coprime-path-nine-2026-10-18
checker=src/coprime_path_structure_certificates.py
python3 -I -S -O "$checker" cuts "$cert_dir/cuts.json.gz"
python3 -I -S -O "$checker" faults "$cert_dir/fault-families.json.gz"
python3 -I -S -O "$checker" extensions "$cert_dir/extensions.json.gz"
\end{verbatim}
\end{minipage}

Require a successful exit, the counts above, and the intended
task parameters: $n=17$, target length nine, minimum cut
side five, and vertex- and edge-deletion budgets of four.
For extensions, require \texttt{side\_size=4} and
\texttt{minimum\_side\_edges=5}. The symmetry checks concern
the host only; no assumption is made that a coloring
respects those symmetries.

\subsection{Lower colorings and the separate SAT certificate}

The complete lower colorings are checked independently by
subset dynamic programming. For each vertex subset and
color, it records the possible endpoints of a path using
exactly that subset. The longest monochromatic path lengths,
measured in edges, are $(8,8)$ at $n=16$ and $(6,6)$ at
$n=12$.

For the $n=13$ upper bound, a separate formula checker verifies
that every non-tautological clause has uniform sign and
consists of the eight distinct edges of a simple path.
The unsatisfiability proof is checked with DRAT-trim, pinned
at revision
{\small\texttt{2e3b2dc0ecf938addbd779d42877b6ed69d9a985}}.
The archive README gives the complete commands.
Its verification record reports a proof core with
57\,506 clauses, 391\,234 lemmas, and 4\,858\,014 resolution
steps, with no RAT lemmas in the core.

These checks have separate roles: path and coloring checkers
validate explicit combinatorial witnesses; the formula
checker validates the path constraints; DRAT-trim validates
the proof of their inconsistency. The $n=17$ argument uses
the explicit host certificates and the written extremal
proof, with no SAT unsatisfiability certificate.

The verification depends on the checkers, their execution
environments, and the published mathematical inputs.
Checking the finite claims does not formally verify the
whole written argument, for which there is no Lean
formalization. Unfinished or time-limited searches supply
no proof of the statements here.

\subsection{Historical records}

The earlier code archive is available at
\href{https://doi.org/10.5281/zenodo.23433640}
{doi:10.5281/zenodo.23433640}.
The earlier manuscript is available at
\href{https://doi.org/10.5281/zenodo.19058647}
{doi:10.5281/zenodo.19058647}.

Both deposits are retained as historical records.
The new nine-edge-path evidence is identified by the
immutable repository revision above, rather than by
these earlier deposits.

\section*{AI assistance and draft status}

Codex assisted with the proof investigation, certificate
software, literature checks, and preparation and editorial
revision of this manuscript. Human verification of every
proof step and citation, and the author's final revision
of the AI-assisted proofs in the author's own words,
remain pending. This is a working draft and has not been
submitted to a journal.

\clearpage



\begin{thebibliography}{99}
\bibitem{ali}
A. A. Ali and W. Staton.
On extremal graphs with no long paths.
\emph{Electronic Journal of Combinatorics}, 3(1):R20, 1996.
\href{https://doi.org/10.37236/1244}{doi:10.37236/1244}.

\bibitem{bgls}
P. N. Balister, E. Gy\H{o}ri, J. Lehel, and R. H. Schelp.
Connected graphs without long paths.
\emph{Discrete Mathematics}, 308(19):4487--4494, 2008.
\href{https://doi.org/10.1016/j.disc.2007.08.047}{doi:10.1016/j.disc.2007.08.047}.
\href{https://www.memphis.edu/msci/people/pbalistr/nopk.pdf}{Author manuscript}.

\bibitem{closure}
J. A. Bondy and V. Chv\'atal.
A method in graph theory.
\emph{Discrete Mathematics}, 15:111--135, 1976.
\href{https://doi.org/10.1016/0012-365X(76)90078-9}{doi:10.1016/0012-365X(76)90078-9}.

\bibitem{du}
Z. Du, W. Xi, Z. Deng, and L. Ma.
Prime certificates for exact vertex-coprime Ramsey numbers.
Preprint, 2026.
\href{https://arxiv.org/abs/2605.26815v2}{arXiv:2605.26815v2}.

\bibitem{eg}
P. Erd\H{o}s and T. Gallai.
On maximal paths and circuits of graphs.
\emph{Acta Mathematica Academiae Scientiarum Hungaricae},
10:337--356, 1959.
\href{https://doi.org/10.1007/BF02024498}{doi:10.1007/BF02024498}.
\href{https://web.vu.lt/mif/s.jukna/tropical/Erdos-Gallai-1959.pdf}{Full text}.

\bibitem{fkv}
Z. F\"uredi, A. Kostochka, and J. Verstra\"ete.
Stability in the Erd\H{o}s--Gallai Theorems on cycles and paths.
\emph{Journal of Combinatorial Theory, Series B}, 121:197--228, 2016.
\href{https://doi.org/10.1016/j.jctb.2016.06.004}{doi:10.1016/j.jctb.2016.06.004}.
Theorem numbering follows
\href{https://arxiv.org/abs/1507.05338v2}{arXiv:1507.05338v2}.

\bibitem{fkvl}
Z. F\"uredi, A. Kostochka, R. Luo, and J. Verstra\"ete.
Stability in the Erd\H{o}s--Gallai Theorem on cycles and paths, II.
\emph{Discrete Mathematics}, 341(5):1253--1263, 2018.
\href{https://doi.org/10.1016/j.disc.2017.12.018}{doi:10.1016/j.disc.2017.12.018}.
Theorem numbering follows
\href{https://arxiv.org/abs/1704.02866v1}{arXiv:1704.02866v1}.

\bibitem{ma-ning}
J. Ma and B. Ning.
Stability results on the circumference of a graph.
\emph{Combinatorica}, 40:105--147, 2020.
\href{https://doi.org/10.1007/s00493-019-3843-4}{doi:10.1007/s00493-019-3843-4}.
The comparison uses
\href{https://arxiv.org/abs/1708.00704v2}{arXiv:1708.00704v2}.

\bibitem{repository}
A. R. Towell.
Computational explorations: coprime path proof notes and certificates.
GitHub repository, revision \texttt{8bda81e}, October 18, 2026.
\href{\repo/tree/\revision}{Immutable source and evidence}.
\end{thebibliography}
\end{document}
